% Shared macros for both pdf and web versions of the blueprint.
% Note: the style guide recommends against \begin{proposition}, but this blueprint uses it throughout for paper fidelity.


% ── Theorem-like environments (these define the dependency graph nodes) ──
\newtheorem{theorem}{Theorem}[chapter]
\newtheorem{proposition}[theorem]{Proposition}
\newtheorem{lemma}[theorem]{Lemma}
\newtheorem{corollary}[theorem]{Corollary}

\theoremstyle{definition}
\newtheorem{definition}[theorem]{Definition}

\theoremstyle{remark}
\newtheorem{remark}[theorem]{Remark}

% ── Mathematical notation ──
\newcommand{\C}{\mathbb{C}}
\newcommand{\F}{\mathbb{F}}
\newcommand{\N}{\mathbb{N}}
\newcommand{\R}{\mathbb{R}}
\newcommand{\E}{\mathbb{E}}
\newcommand{\Prb}{\mathbb{P}}
\newcommand{\id}{\mathrm{Id}}
\newcommand{\Id}{\mathrm{Id}}
\newcommand{\tr}{\operatorname{tr}}
\DeclareMathOperator{\Tr}{Tr}
\newcommand{\Mat}{\operatorname{Mat}}
\newcommand{\code}{\mathcal{C}}
\newcommand{\hilb}{\mathcal{H}}
\newcommand{\algebra}{\mathcal{A}}
\newcommand{\eps}{\varepsilon}
\newcommand{\ot}{\otimes}
\newcommand{\ket}[1]{\lvert #1 \rangle}
\newcommand{\bra}[1]{\langle #1 \rvert}
\newcommand{\braket}[2]{\langle #1 \mid #2 \rangle}
\newcommand{\bone}[1]{\mathbf{1}\!\left[#1\right]}
\newcommand{\polyfunc}[3]{\mathcal{P}(#1,#2,#3)}
\newcommand{\polysub}[3]{\mathrm{PolySub}(#1,#2,#3)}
\newcommand{\polymeas}[3]{\mathrm{PolyMeas}(#1,#2,#3)}
\DeclareMathOperator{\poly}{poly}
\newcommand{\calA}{\mathcal{A}}
\newcommand{\calB}{\mathcal{B}}
\newcommand{\calD}{\mathcal{D}}
\newcommand{\calH}{\mathcal{H}}
\newcommand{\calK}{\mathcal{K}}
\newcommand{\calX}{\mathcal{X}}
\newcommand{\abs}[1]{\lvert #1 \rvert}

% ── Bold random-variable notation (from the paper's wright.sty) ──
\newcommand{\ba}{\boldsymbol{a}}
\newcommand{\bg}{\boldsymbol{g}}
\newcommand{\bh}{\boldsymbol{h}}
\newcommand{\bi}{\boldsymbol{i}}
\newcommand{\br}{\boldsymbol{r}}
\newcommand{\bs}{\boldsymbol{s}}
\newcommand{\bu}{\boldsymbol{u}}
\newcommand{\bv}{\boldsymbol{v}}
\newcommand{\bw}{\boldsymbol{w}}
\newcommand{\bx}{{\boldsymbol{x}}}
\newcommand{\by}{\boldsymbol{y}}
\newcommand{\bz}{\boldsymbol{z}}
